\documentclass[11pt, a4paper]{article}

\usepackage[utf8]{inputenc}
\usepackage[T1]{fontenc}
\usepackage[french]{babel}
\usepackage{amsmath}
\usepackage{amssymb}
\usepackage{bm}
\usepackage{geometry}
\usepackage{booktabs}
\usepackage{tcolorbox}
\geometry{margin=2.5cm}

\title{\textbf{Ferroélectricité \\ Démonstration variationnelle des équations de champ}}
\author{Note de Synthèse - Modélisation Couplée}
\date{}

\begin{document}

\maketitle

\begin{abstract}
Dérivation rigoureuse des équations de Gauss, d'équilibre mécanique et de Ginzburg-Landau par principe variationnel appliqué à la fonctionnelle de Landau-Ginzburg-Devonshire (LGD).
\end{abstract}

\section{Notation et conventions}

Convention de sommation d'Einstein sur les indices répétés $i,j,k,\ldots \in \{1,2,3\}$. Notation de dérivée partielle : $\partial_i = \partial/\partial x_i$ et $f_{,i} := \partial_i f$. Le domaine $\Omega \subset \mathbb{R}^3$ possède une frontière $\partial\Omega$ régulière de normale extérieure unitaire $n_i$.

\paragraph{Symétries des tenseurs constitutifs :}
\begin{equation*}
c_{ijkl} = c_{jikl} = c_{ijlk} = c_{klij}, \qquad Q_{ijkl} = Q_{jikl} = Q_{ijlk}, \qquad G_{ijkl} = G_{klij}
\end{equation*}

\section{Variables d'état et cinématique}

\subsection{Champs primaires}
\begin{table}[h]
\centering
\begin{tabular}{ll}
\toprule
\textbf{Champ} & \textbf{Définition} \\
\midrule
Déplacement & $\mathbf{u}$ : champ mécanique vectoriel $u_i : \Omega \to \mathbb{R}$ \\
Polarisation & $\mathbf{P}$ : paramètre d'ordre ferroélectrique $P_i : \Omega \to \mathbb{R}$ \\
Potentiel électrique & $\phi$ : champ scalaire électrostatique $\phi : \Omega \to \mathbb{R}$ \\
\bottomrule
\end{tabular}
\end{table}

\subsection{Relations cinématiques fondamentales}

Tenseur de déformation (hypothèse des petites perturbations) :
\begin{equation}
\varepsilon_{ij} = \frac{1}{2}(u_{i,j} + u_{j,i})
\end{equation}

Déformation propre électrostrictive (eigenstrain spontanée induite par la polarisation) :
\begin{equation}
\varepsilon^{(0)}_{ij} = Q_{ijkl}P_k P_l
\end{equation}

Déformation élastique, tenseur des contraintes de Cauchy, et déplacement électrique :
\begin{equation}
\varepsilon^{(e)}_{ij} = \varepsilon_{ij} - Q_{ijkl}P_k P_l, \qquad \sigma_{ij} = c_{ijkl}\varepsilon^{(e)}_{kl}, \qquad D_i = \varepsilon_0\kappa^b E_i + P_i
\end{equation}

De l'irrotationnalité quasi-statique $\nabla \times \mathbf{E} = \mathbf{0}$, il existe $\phi$ tel que $E_i = -\phi_{,i}$, d'où :
\begin{equation}
D_i = -\varepsilon_0\kappa^b\phi_{,i} + P_i
\end{equation}
où $\kappa^b$ est la constante diélectrique de fond (contribution ionique et électronique du réseau).

\section{Fonctionnelle d'énergie libre}

La fonctionnelle de Helmholtz totale est définie sur $\Omega$ par :
\begin{equation}
\mathcal{F}[\mathbf{P},\mathbf{u},\phi] = \int_\Omega \mathcal{H} \, dV, \qquad \mathcal{H} = f_L + f_g + f_e + f_{em}
\end{equation}

\paragraph{Potentiel de Landau-Devonshire (double puits) :}
Développement polynomial de $\mathcal{H}$ en $\mathbf{P}$ respectant la symétrie cubique $m\bar{3}m$ :
\begin{equation}
f_L = \alpha_1\sum_i P_i^2 + \alpha_{11}\sum_i P_i^4 + \alpha_{12}\sum_{i < j}P_i^2 P_j^2 + \alpha_{111}\sum_i P_i^6 + \alpha_{112}\sum_{i\neq j}P_i^4 P_j^2 + \alpha_{123}P_1^2 P_2^2 P_3^2
\end{equation}
Ou sous forme tensorielle compacte : $f_L = A_{ij}P_iP_j + B_{ijkl}P_iP_jP_kP_l + C_{ijklmn}P_iP_jP_kP_lP_mP_n$. \\
Loi de Curie-Weiss : $\alpha_1 = (T - T_c)/(2\varepsilon_0 C)$ avec $T_c$ température de Curie et $C > 0$. Le double puits ($\alpha_1 < 0$) traduit la brisure de symétrie ferroélectrique.

\paragraph{Énergie de gradient de Ginzburg (parois de domaine) :}
\begin{equation}
f_g = \frac{1}{2} G_{ijkl} (\partial_j P_i) (\partial_l P_k)
\end{equation}
Pénalise les variations spatiales de $\mathbf{P}$ : contrôle l'épaisseur et l'énergie des parois de domaine. Cas isotrope : $G_{ijkl} = g\delta_{ik}\delta_{jl} \implies f_g = \frac{g}{2}|\nabla\mathbf{P}|^2$.

\paragraph{Énergie élastique avec couplage électrostrictif :}
\begin{equation}
f_e = \frac{1}{2} c_{ijkl} \left(\varepsilon_{ij} - Q_{ijmn}P_mP_n\right) \left(\varepsilon_{kl} - Q_{klpq}P_pP_q\right)
\end{equation}

\paragraph{Enthalpie électrique (couplage champ-polarisation) :}
\begin{equation}
f_{em} = -\frac{1}{2}\varepsilon_0\kappa^b E_i E_i - P_i E_i = \frac{\varepsilon_0\kappa^b}{2}\phi_{,i}\phi_{,i} + P_i\phi_{,i}
\end{equation}

\section{Équation électrostatique (Théorème I - Loi de Gauss)}

Soit $\delta\phi \in H^1_0(\Omega)$ une fonction-test admissible. La stationnarité de $\mathcal{F}$ impose $\delta_\phi\mathcal{F} = 0$ :
\begin{equation}
\delta_\phi\mathcal{F} = \left.\frac{d}{d\varepsilon}\mathcal{F}[\mathbf{P},\mathbf{u},\phi+\varepsilon\delta\phi]\right|_{\varepsilon=0} = \int_\Omega \frac{\partial f_{em}}{\partial \phi_{,i}} \partial_i(\delta\phi) \, dV
\end{equation}
Avec $\frac{\partial f_{em}}{\partial \phi_{,i}} = \varepsilon_0\kappa^b\phi_{,i} + P_i$. L'intégration par parties donne :
\begin{equation}
\delta_\phi\mathcal{F} = -\int_\Omega \partial_i \left(\varepsilon_0\kappa^b\phi_{,i} + P_i\right)\delta\phi \, dV + \oint_{\partial\Omega} \left(\varepsilon_0\kappa^b\phi_{,i}+P_i\right)n_i \delta\phi \, dS
\end{equation}
Le terme de surface s'annule car $\delta\phi\big|_{\partial\Omega} = 0$. Par le lemme fondamental du calcul des variations :

\begin{tcolorbox}[title=Équation I — Loi de Gauss / Équation de Poisson]
\begin{equation}
\nabla\cdot\mathbf{D} = 0 \qquad\Longleftrightarrow\qquad \varepsilon_0\kappa^b\Delta\phi = \nabla\cdot\mathbf{P}
\end{equation}
\end{tcolorbox}

\section{Équilibre mécanique (Théorème II - Équation de Newton)}

Soit $\delta u_i \in H^1_0(\Omega)$. On identifie d'abord le tenseur des contraintes :
\begin{equation}
\sigma_{ij} := \frac{\partial f_e}{\partial \varepsilon_{ij}} = c_{ijkl}\left(\varepsilon_{kl} - Q_{klmn}P_mP_n\right) = c_{ijkl}\varepsilon^{(e)}_{kl}
\end{equation}
\begin{equation}
\delta_u\mathcal{F} = \int_\Omega \sigma_{ij}\delta\varepsilon_{ij} \, dV = -\int_\Omega (\partial_j\sigma_{ij})\delta u_i \, dV + \oint_{\partial\Omega}\sigma_{ij}n_j\delta u_i \, dS = 0
\end{equation}
Ce qui donne l'équation d'équilibre local :

\begin{tcolorbox}[title=Équation II — Équilibre mécanique de Lamé-Navier généralisé]
\begin{equation}
\partial_j\sigma_{ij} = 0 \qquad\Longleftrightarrow\qquad \nabla\cdot\boldsymbol{\sigma} = \mathbf{0}
\end{equation}
\end{tcolorbox}

\section{Équation de champ de phase (Théorème III - Équation de Ginzburg-Landau)}

La dérivée variationnelle de $\mathcal{F}$ par rapport à $P_k$ s'écrit :
\begin{equation}
\frac{\delta\mathcal{F}}{\delta P_k} = \frac{\partial\mathcal{H}}{\partial P_k} - \partial_j\left(\frac{\partial\mathcal{H}}{\partial(\partial_j P_k)}\right)
\end{equation}
Calcul des contributions :
\begin{align}
\text{Landau :} &\quad \frac{\partial f_L}{\partial P_k} = 2A_{kj}P_j + 4B_{kjlm}P_jP_lP_m + 6C_{kjlmnp}P_jP_lP_mP_nP_p \\
\text{Ginzburg :} &\quad -\partial_j\left(\frac{\partial f_g}{\partial(\partial_j P_k)}\right) = -G_{kjlm}\partial_j\partial_l P_m \\
\text{Piézo inverse :} &\quad \frac{\partial f_e}{\partial P_k} = -2c_{ijlm}\varepsilon^{(e)}_{ij}Q_{lmkq}P_q = -2Q_{ijkl}P_l\sigma_{ij} \\
\text{Électrique :} &\quad \frac{\partial f_{em}}{\partial P_k} = \phi_{,k} = -E_k
\end{align}
L'équation de Ginzburg-Landau dépendante du temps (TDGL) postule $\Gamma\frac{\partial P_k}{\partial t} = -\frac{\delta\mathcal{F}}{\delta P_k}$ :

\begin{tcolorbox}[title=Équation III — Champ de phase / TDGL (Allen-Cahn ferroélectrique)]
\begin{equation}
\Gamma\frac{\partial P_k}{\partial t} = -\frac{\partial f_L}{\partial P_k} + G_{kjlm}\partial_j\partial_l P_m + 2Q_{ijkl}P_l\sigma_{ij} + E_k
\end{equation}
\end{tcolorbox}

\section{Système couplé et Conditions aux limites}

Le système repose sur la résolution simultanée de trois EDP non-linéaires fortement couplées :
\begin{align}
\varepsilon_0\kappa^b\Delta\phi &= \nabla\cdot\mathbf{P} \\
\partial_j\left[c_{ijkl}\left(\varepsilon_{kl} - Q_{klmn}P_mP_n\right)\right] &= 0 \\
\Gamma\dot{P}_k &= -\frac{\partial f_L}{\partial P_k} + G_{kjlm}P_{m,jl} + 2Q_{ijkl}P_l\sigma_{ij} + E_k
\end{align}

\paragraph{Conditions de bord standard :}
\begin{table}[h]
\centering
\begin{tabular}{lp{10cm}}
\toprule
\textbf{Champ} & \textbf{Condition sur $\partial\Omega$} \\
\midrule
Mécanique & $\sigma_{ij}n_j = t_i$ (traction) ou $u_i = \bar{u}_i$ (déplacement imposé) \\
Électrostatique & $\phi = \bar{\phi}$ (tension) ou $D_i n_i = \sigma_f$ (charge libre de surface) \\
Champ de phase & $G_{kjlm}(\partial_l P_m)n_j = 0$ (condition de Neumann naturelle, paroi neutre) \\
\bottomrule
\end{tabular}
\end{table}

\end{document}